% GENERATED FILE. Edit canonical lessons, then run npm run generate:books.
% canonical-source-hash: fnv1a64:fbf794a78b46589e
% canonical-lessons: HI-C558-khvahis, HI-C558-divana, HI-C558-ruci, HI-C558-akanksa, HI-C558-junun

\chapter{Longings and Passions}
\label{ch:hi-longings-and-passions}
\hlchaptermodality{558}

\begin{chapteropening}
\textbf{By the end of this chapter:} \emph{I can talk about wishes, ambitions and passions.}

Complete the last lesson of chapter 558: I can talk about wishes, ambitions and passions.
\end{chapteropening}

\section[\texorpdfstring{\hi{ख़्वाहिश} (khvāhiś)}{ख़्वाहिश (khvāhiś)}]{\hi{ख़्वाहिश} (khvāhiś) — a wish, a longing (for something)}

\label{lesson:HI-C558-khvahis}



\begin{quote}
Before the new one: say the Hindi for misfortune, then the Hindi for fear, dread.
\end{quote}

\subsection*{You'll want to know: \hi{ख़्वाहिश}}
\textbf{\hi{ख़्वाहिश}} — \emph{khvāhiś} — \textquotedblleft{}a wish, a longing (for something)\textquotedblright{}.

\begin{quote}
\hi{दादी} \hi{ने} \hi{मेरी} \hi{हर} \hi{ख़्वाहिश} \hi{पूरी} \hi{की।}
\end{quote}

\begin{quote}
\emph{Dādī ne merī har khvāhiś pūrī kī.}
\end{quote}

\begin{quote}
Grandmother granted my every wish.
\end{quote}

Feminine.

\begin{cousinweb}
\textbf{\hi{ख़्वाहिश}} (\emph{khvāhiś}) is usually traced to Persian \emph{khwāhish}, \textquotedblleft{}a wish\textquotedblright{}.
\end{cousinweb}

\begin{grammarlens}[title={What you've built}]
One more word for saying what happens if.
\end{grammarlens}

\subsection*{Your turn}
\begin{itemize}
\raggedright
  \item \emph{Say it:} \emph{khvāhiś}
  \item \emph{Say it:} \emph{khvāhiś}, once more
  \item \emph{Say it:} \emph{Dādī ne merī har khvāhiś pūrī kī.} — Grandmother granted my every wish.
  \item \emph{Recall:} say the Hindi for misfortune, then the Hindi for fear, dread, then say \emph{khvāhiś} again
\end{itemize}

\begin{tcolorbox}[breakable,colback=teal!4,colframe=teal!35!black,title={Before you move on}]
What does \hi{ख़्वाहिश} mean? (\textquotedblleft{}A wish, a longing (for something)\textquotedblright{}.) Say it once more.
\end{tcolorbox}
\section[\texorpdfstring{\hi{दीवाना} (dīvānā)}{दीवाना (dīvānā)}]{\hi{दीवाना} (dīvānā) — crazy (about); mad (for)}

\label{lesson:HI-C558-divana}



\begin{quote}
Before the new one: say the Hindi for fear, dread, then the Hindi for a wish, a longing (for something).
\end{quote}

\subsection*{You'll want to know: \hi{दीवाना}}
\textbf{\hi{दीवाना}} — \emph{dīvānā} — \textquotedblleft{}crazy (about); mad (for)\textquotedblright{}.

\begin{quote}
\hi{मीरा} \hi{पुराने} \hi{गानों} \hi{की} \hi{दीवानी} \hi{है।}
\end{quote}

\begin{quote}
\emph{Mīrā purāne gānoṁ kī dīvānī hai.}
\end{quote}

\begin{quote}
Mira is crazy about old songs.
\end{quote}

Agrees like an adjective: \hi{दीवाना}, \hi{दीवानी}, \hi{दीवाने}.

\begin{cousinweb}
\textbf{\hi{दीवाना}} (\emph{dīvānā}) is usually traced to Persian \emph{dīvāna}, \textquotedblleft{}mad, possessed\textquotedblright{}; Hindi uses it for someone wild about a thing or a person.
\end{cousinweb}

\begin{grammarlens}[title={What you've built}]
One more word for saying what happens if.
\end{grammarlens}

\subsection*{Your turn}
\begin{itemize}
\raggedright
  \item \emph{Say it:} \emph{dīvānā}
  \item \emph{Say it:} \emph{dīvānā}, once more
  \item \emph{Say it:} \emph{Mīrā purāne gānoṁ kī dīvānī hai.} — Mira is crazy about old songs.
  \item \emph{Recall:} say the Hindi for fear, dread, then the Hindi for a wish, a longing (for something), then say \emph{dīvānā} again
\end{itemize}

\begin{tcolorbox}[breakable,colback=teal!4,colframe=teal!35!black,title={Before you move on}]
What does \hi{दीवाना} mean? (\textquotedblleft{}Crazy (about); mad (for)\textquotedblright{}.) Say it once more.
\end{tcolorbox}
\section[\texorpdfstring{\hi{रुचि} (ruci)}{रुचि (ruci)}]{\hi{रुचि} (ruci) — an interest; a liking (for something)}

\label{lesson:HI-C558-ruci}



\begin{quote}
Before the new one: say the Hindi for a wish, a longing (for something), then the Hindi for crazy (about).
\end{quote}

\subsection*{You'll want to know: \hi{रुचि}}
\textbf{\hi{रुचि}} — \emph{ruci} — \textquotedblleft{}an interest; a liking (for something)\textquotedblright{}.

\begin{quote}
\hi{मीरा} \hi{की} \hi{संगीत} \hi{में} \hi{गहरी} \hi{रुचि} \hi{है।}
\end{quote}

\begin{quote}
\emph{Mīrā kī saṅgīt meṁ gahrī ruci hai.}
\end{quote}

\begin{quote}
Mira has a deep interest in music.
\end{quote}

Feminine.

\begin{cousinweb}
\textbf{\hi{रुचि}} (\emph{ruci}) is a tatsama, taken whole from Sanskrit \emph{ruci}, \textquotedblleft{}light; liking, taste, relish\textquotedblright{} (Monier-Williams).
\end{cousinweb}

\begin{grammarlens}[title={What you've built}]
One more word for saying what happens if.
\end{grammarlens}

\subsection*{Your turn}
\begin{itemize}
\raggedright
  \item \emph{Say it:} \emph{ruci}
  \item \emph{Say it:} \emph{ruci}, once more
  \item \emph{Say it:} \emph{Mīrā kī saṅgīt meṁ gahrī ruci hai.} — Mira has a deep interest in music.
  \item \emph{Recall:} say the Hindi for a wish, a longing (for something), then the Hindi for crazy (about), then say \emph{ruci} again
\end{itemize}

\begin{tcolorbox}[breakable,colback=teal!4,colframe=teal!35!black,title={Before you move on}]
What does \hi{रुचि} mean? (\textquotedblleft{}An interest; a liking (for something)\textquotedblright{}.) Say it once more.
\end{tcolorbox}
\section[\texorpdfstring{\hi{आकांक्षा} (ākāṅkṣā)}{आकांक्षा (ākāṅkṣā)}]{\hi{आकांक्षा} (ākāṅkṣā) — an ambition; an aspiration}

\label{lesson:HI-C558-akanksa}



\begin{quote}
Before the new one: say the Hindi for crazy (about), then the Hindi for an interest.
\end{quote}

\subsection*{You'll want to know: \hi{आकांक्षा}}
\textbf{\hi{आकांक्षा}} — \emph{ākāṅkṣā} — \textquotedblleft{}an ambition; an aspiration\textquotedblright{}.

\begin{quote}
\hi{उसकी} \hi{आकांक्षा} \hi{एक} \hi{बड़ा} \hi{लेखक} \hi{बनने} \hi{की} \hi{है।}
\end{quote}

\begin{quote}
\emph{Uskī ākāṅkṣā ek baṛā lekhak banne kī hai.}
\end{quote}

\begin{quote}
His ambition is to become a great writer.
\end{quote}

Feminine.

\begin{cousinweb}
\textbf{\hi{आकांक्षा}} (\emph{ākāṅkṣā}) is a tatsama, taken whole from Sanskrit \emph{ākāṅkṣā}, \textquotedblleft{}desire, wish\textquotedblright{} (Monier-Williams).
\end{cousinweb}

\begin{grammarlens}[title={What you've built}]
One more word for saying what happens if.
\end{grammarlens}

\subsection*{Your turn}
\begin{itemize}
\raggedright
  \item \emph{Say it:} \emph{ākāṅkṣā}
  \item \emph{Say it:} \emph{ākāṅkṣā}, once more
  \item \emph{Say it:} \emph{Uskī ākāṅkṣā ek baṛā lekhak banne kī hai.} — His ambition is to become a great writer.
  \item \emph{Recall:} say the Hindi for crazy (about), then the Hindi for an interest, then say \emph{ākāṅkṣā} again
\end{itemize}

\begin{tcolorbox}[breakable,colback=teal!4,colframe=teal!35!black,title={Before you move on}]
What does \hi{आकांक्षा} mean? (\textquotedblleft{}An ambition; an aspiration\textquotedblright{}.) Say it once more.
\end{tcolorbox}
\section[\texorpdfstring{\hi{जुनून} (junūn)}{जुनून (junūn)}]{\hi{जुनून} (junūn) — a passion; an obsession}

\label{lesson:HI-C558-junun}



\begin{quote}
Before the new one: say the Hindi for an interest, then the Hindi for an ambition.
\end{quote}

\subsection*{You'll want to know: \hi{जुनून}}
\textbf{\hi{जुनून}} — \emph{junūn} — \textquotedblleft{}a passion; an obsession\textquotedblright{}.

\begin{quote}
\hi{उसे} \hi{क्रिकेट} \hi{का} \hi{जुनून} \hi{है।}
\end{quote}

\begin{quote}
\emph{Use krikeṭ kā junūn hai.}
\end{quote}

\begin{quote}
He is crazy about cricket.
\end{quote}

Masculine.

\begin{cousinweb}
\textbf{\hi{जुनून}} (\emph{junūn}) is usually traced to Arabic \emph{junūn}, \textquotedblleft{}madness, possession\textquotedblright{}, through Persian. Hindi uses it for a consuming passion.
\end{cousinweb}

\begin{grammarlens}[title={What you've built}]
One more word for saying what happens if.
\end{grammarlens}

\subsection*{Your turn}
\begin{itemize}
\raggedright
  \item \emph{Say it:} \emph{junūn}
  \item \emph{Say it:} \emph{junūn}, once more
  \item \emph{Say it:} \emph{Use krikeṭ kā junūn hai.} — He is crazy about cricket.
  \item \emph{Recall:} say the Hindi for an interest, then the Hindi for an ambition, then say \emph{junūn} again
\end{itemize}

\begin{tcolorbox}[breakable,colback=teal!4,colframe=teal!35!black,title={Before you move on}]
What does \hi{जुनून} mean? (\textquotedblleft{}A passion; an obsession\textquotedblright{}.) Say it once more.
\end{tcolorbox}
